% Note that because I talk a lot about binding ... this discussion within my methods section may be quite important for examing members that understands the physics but do not understand how MD simulations in particular is rigoroous enough to make any of those predictions.  Having those details here means that I understand how that link is established and that I'm not just making up stories based on pretty pictures.  These are some of the goals that I want to acheive for the benefit of the reader of this chapter.


% From D. Caplan's thesis - Molecular mechanics (MM) is the empirical mathematical representation of atomic inter- actions in the classical limit. MM was developed over 50 years ago and has been applied to organic chemistry as a tool to estimate specific energetic properties of small molecules. Since then it has evolved into various functional forms (known as �force fields�) used to describe the potential energy of a system by computing atomic interactions. Force fields may vary in the way that they represent atomic particles to balance experimental accu- racy with computational efficiency. Molecular dynamics (MD) is a MM-based method in which force fields are used to calculate the position, velocity, and acceleration of atoms over time to produce a trajectory.

\chapter[Methods]{Methods: Molecular Modelling and Simulation}

\section{Molecular Mechanics}
% Start by describing quantum mechanics 
% Assumptions and limitations of molecular mechanics
In an ideal world, one would be able to predict the evolution of a molecular system by solving the Schrodinger's equation, which predicts the t-evolution of all the atomic-nuclei and electronic motion.  However, to do so for a biomolecular system is extremely computationally-prohibitive, even for a very small molecule. REF

In the past decades, molecular mechanics, a classical mechanics approach, was developed to model and simulate a large system possible. In molecular mechanics, an empirical potential energy function which depends only on the positions of atoms, is employed to describe the molecular energies of molecules. Simply put, molecules are treated as a set of masses connected by springs.

The underlying theoretical basis for molecular mechanics is rooted in the Born-Oppenheimer approximation. Under this approximation, electronic and nuclear motions are considered to be uncoupled, and therefore are treated separately. For this reason, the system's total energy (the Hamiltonian, $H$) only depends on the electronic motion, and the nuclei are assumed to be in fixed positions. 
% This is because the nuclei, which are much heavier than electrons, are typically fixed on the time scales of electronic vibration.

This potential energy function employed in molecular mechanics corresponds to the Born-Oppenheimer energy surface of the molecular system in its electronic ground-state. The movements of all atomic nuclei can be treated as a point charge in space, where the movement of electrons is implicitly accounted for in the specific form of the PEF and its associated parameter set, or force field.  Because electronically excited states are not treated, no chemical reactions occur in the system. However, although MM tells us nothing about electronic characteristics, its computational speed and its ability to deal with large systems makes it very attractive.
% \cite{M. Born and R. Oppenheimer, Quantum Theory of Molecules, Ann. Physics, vol. 84, pp. 457�484, 1927.} -- can't find this article.

\section{Force Fields}

In the past four decades, a variety of functional forms have been developed, where each is optimized for a specific purpose have been developed for the purpose of modelling and simulating biomolecular systems. The functional form together with its associated parameters defines a \emph{force field}.

A typical empirical force field has many parameters to be determined.  These parameters are typically determined by fitting to quantum mechanical calculations, and to physical measurements (e.g., using data from X-ray crystallography, neutron diffraction, vibrational raman spectroscopy) of small organic compounds or building blocks of peptides and nucleic acids.  Force fields are iteratively improved and extended by predicting experimentally observable quantities for small compounds, and adjusting parameters by comparing predicted quantities with experimental measurements.

Currently, many force fields for biopolymers such as peptides, proteins, nucleic acids and carbohydrates have been developed. Examples of these are the force fields AMBER, CHARMM, GROMOS, and OPLS. REFs  Although each force field have their variation of the PEF form, these force fields in general adopt the following general form:

  \begin{equation}
    \label{eq:ff_general}
    V = V(bonds) + V(angles) + V(impropers) + V(dihedrals) + V(non-bonded pair interactions)
  \end{equation}

The first four terms of \ref{eq:ff_general}, $V(bonds)$, $V(angles)$, and $V(improper)$  are the functional terms used to describe bonded interactions and dynamics within a molecule.
These motions consists of covalent bond stretching, angle-bending, and torsional rotations.  It is important to note that bonds cannot be formed nor broken. 

Covalent bond stretching can be described using Hooke's Law for harmonic motion,
\begin{equation}
 V(bonds) = \sum_{bonds} k_b(b - b_0)^2.
\end{equation}

Similarly, angle-bending motions can be describing using a similar potential function,
\begin{equation}
V(angles) = \sum_{angles} k_{\theta}(\theta - \theta_{0})^2, \\
\end{equation}
where the constants $k_{b}$ and $k_{\theta}$ are the equilibrium force constants for bonds and angles, respectively, and $b_{0}$ and $\theta_{0}$ are the reference bond length and angle, respectively. Note that total energetic contribution from bond stretching and angle bending motions are computed by summing over all bonds and angles in the system.

Non-bonded interaction potentials predominantly account for van der Waals and electrostatic forces.  These are important for the transferability of a force field, that is, the ability to reuse parameters created based on one molecule to predict properties of other molecules.  The non-bonded forces are only applied to atom pairs separated by at least three bonds. Van der Waals forces are represented by the Lennard-Jones potential,

\begin{equation}
  \label{eq:lj_function}
  V_{LJ}(r) = (\frac{r_{min}}{r})^{12} - 2(\frac{r_{min}}{r})^6
\end{equation}

The first term of \ref{eq:lj_function} is a repulsive force, and the second term is an attractive force at atomic distances where atomic shells overlap.  Here the constant $r_{min}$ is the preferred interaction distance for the atom pair or the van der Waals radii of atoms. Together these constitute a rapidly decaying function which represents a short-range attractive interaction that is characteristic of van der Waals forces. Parameterization of van der Waals potential to accurately reproduce these types of intermolecular interactions are typically challenging to determine because of the need for accurate reproduction of atomic radii, weak atomic attraction, and repulsive forces.
% Transferability (i.e., the ability to reuse parameters created based on one molecule to predict properties of another molecule 

Furthermore, the reproduction of electrostatics is paramount for describing interactions of polar and charged groups in biomolecules. To do so, an additive pairwise Coulombic potential is applied between all atom pairs which carry partial charges or charges at their nuclear centres.  The Coulombic potential, which depends on $r$, the distance between the pair of atoms, is given by

\begin{equation}
  \label{eq:electrostatic}
  V_{coulombic}(r) = \frac{q_iq_j}{\epsilon r}, 
\end{equation}
where $q_{i}$ and $q_{j}$ are the point charges on atoms $i$ and $j$, respectively, and $\epsilon$ is the the permittivity of vacuum.  In contrast to van der Waals interactions,  electrostatic interactions are long-range interactions. Accordingly depicted in Figure XXX,  in comparisons to the Lennard-Jones potential, the Coulombic potential is a much slower decaying function, indicating that electrostatic effects are significant even up to distances of 1 - 1.5 nm.

\begin{figure}
 \centering
 \includegraphics[width=5in]{figures/methods/LJ_Coulumb_potentials.pdf}
 \caption[Lennard-Jones and Coulombic Potentials]{Lennard-Jones and Coulombic Potentials.}
 \label{fig:nonbonded_potentials}
\end{figure}

As previously discussed in Chapter 1 \ref{chap:introduction}, hydrogen bonding is an important type of electrostatic interaction and hence their proper representation must be captured by a force field.
Hydrogen bonding interactions are generally considered to be accounted for by the combination of van der Waals and electrostatic potentials.  However, an explicit term solely for describing hydrogen bonding may be employed by some force fields.  For example, the original AMBER force field introduced in 1984 added a function of the form $\sum_{hbonds} C_{ij}/R_{ij}^{12} - D_{ij}/R_{ij}^{10}$ to account for hydrogen bonds.\cite{Weiner:1984uw}


% Other things that I will not be discussing -- what were they again?

\subsection{Electronic Polarizability}
Polarizability is the measure of the change in a molecule's electron distribution in response to an applied electric field, and is a property of matter which can be induced by electrostatic interactions with solvents or ionic reagents. 

As discussed above, most current force fields use a fixed-charge model where each atom in the system is assigned a single value for the atomic charge.  This charge remains independent of changes in its local electrostatic environment. Next-generation force fields have been proposed to  incorporate models for electronic polarizability. REF However, because of the prohibitive computational cost associated with the use of these force fields, they have not been widely adopted by modern simulations.


% The introduction of polarizable force fields into common use has been inhibited by the high computational cost associated with solving this problem self-consistently.


\subsection{The OPLS Force Field}

I need to read a crap load about the OPLS-AA force field !!!
The general form of the force field potential energy function used in OPLS-AA is
  
  \begin{equation}
    \begin{split}
          E = \sum_{bonds} k_b(b-b_0)^2 
          + \sum_{angles} k_{\theta}(\theta - \theta_{0})^2 \\
          + \sum_{dihedrals} k_{\chi}(1 + cos(n\chi - \delta)) 
          + \sum_{impropers} k_{\gamma}(\phi - \phi_{0})^2 \\
          + \sum_{nonbonded} \frac{q_1q_2}{er} \\
          + \sum_{nonbonded} \epsilon [(\frac{r_{min}}{r})^{12} - 2(\frac{r_{min}}{r})^6]
    \end{split}
  \end{equation}


\section{Molecular Dynamics Simulation}

MD is a molecular mechanics simulation technique.

% I like these analogies and I think it will help make my methods sections less dry
In many respects, Molecular Dynamics simulations are very similar to real experiments. When we perform a real experiment, we first prepare a sample of the material that we wish to study. This sample is then connected to a measuring instrument from which we can measure the property of interest during a certain time interval. If our measurements are subject to statistical noise (as most measurements are), then the longer we average, the more accurate our measurement becomes. 

The same approach is applied when performing a Molecular Dynamics simulation. First, a model system consisting of N particles is selected and we record the dynamics of this system (via solving Newton's equation of motion) until the properties of the system no longer change with time (we equilibrate the system).  After equilibration, we perform the actual measurement. In fact, some of the most common mistakes that can be made when performing a computer experiment are very similar to the mistakes that can be made in real experiments (e.g., the sample is not prepared correctly, the measurement is too short, the system undergoes an irreversible change during the experiment, or we do not measure what we think). (Copied and adapted from Molecular Simulations book)

MD simulations employ an empirical mathematical function to describe the atomic interactions in a molecular system, and together with classical laws of Newtonian mechanics, to simulate systems undergoing thermal motion from which atomic trajectories are generated. 

Thermodynamic and kinetic properties can then be extracted as time averages from these trajectories and used to make a number of predictions that are often experimentally challenging to observe or measure.

MD simulations have been used in recent years to gain insight into many fundamental problems in biology and biochemistry, including protein dynamics and function, protein folding, self-aggregation, and protein-ligand binding. This classical mechanics approach has been shown to be able to provide a valid approximation of many biological systems thus far. REFs

% Approximations made in MD.
% Note the audience of your thesis.  It is good to cover basic ground.  I think I should also brief talk about enhanced techniques.

\subsection{Core Algorithm}
% This is the part with the time propagation using the Equations of motion

The outcome of MD simulations is controlled by a force field, which as discussed above, is a potential energy function describing both intra and intermolecular interactions, which only depends on the positions of the atoms in the system.

For a system of $N$ interacting atoms, the force on each atom pair in the system is determined by the spatial derivative of the potential energy of the atom pair.  % Not exactly ... according to Chris Madill's thesis, U is the potential energy between an atom pair
\begin{equation}
% F_i = - \nabla U
  F_i = - \frac{\delta V}{\delta r_i}, i = 1, 2, 3 ,..., N
\end{equation}

where $V$ is the potential energy of the system. The calculated force vectors are summed together, yielding the net force vector acting on every atom in the system. 

In MD simulations, $V$ is approximated using a molecular mechanics force field, a set of parameters which describes the interaction energies of the atoms in the system.

The acceleration $a_i$ of each atom is given by,
\begin{equation}
a_i = \frac{F_i}{m_i}
\end{equation}
where $m_i$ is the mass of the atom $i$.

The next position of the atom at time $ t + \delta t$
\begin{equation}
x_i(t + \delta t) = x_i(t) + v_i(t)\delta t + \frac{a_i(t)\delta t^2}{2}  
\end{equation}

The above equations are solved numerically using a time integration algorithm using a small timestep. For numerical stability reasons, a small integration step in the range of 1 - 4 femtoseconds is employed. Typically, a timestep of 2 fs is used, which is twice the time period for the fastest vibrational motion (for bonds involving hydrogen), is used in MD simulations of biomolecular systems. % [Ref: Chris Madill's and Tom's thesis]

There are many different forms of integration algorithms, with the most common being the Verlet and leap-frog algorithms. 
% According to gromacs manual:
% md - is a leap-frog algorithm for integrating Newton's equation.
% Note that in absence of pressure and temperature coupling, verlet and leap-frog is equivalent ... and generates identical trajectories of motion.

% How much detail? Talk about the velocity verlet integration algorithm ?
A simulation proceeds iteratively: once all of the new positions for all atoms are predicted, interatomic forces are updated based on these new coordinates, and this entire process is repeated.

% Marty's thesis included a little bit on the verlet algorithm.  CN did not.

\subsection{Temperature and Pressure}
Here describe how temperature and pressure coupling is handled in MD at a high-level without going into the details of the math.

The inclusion of temperature in MD allows for the account of entropy in our systems, which is especially important for free energy determination.  I copied below from ``Molecular simulations the book''

For instance, a convenient definition of the temperature in a (classical) many-body system makes use of the equipartition of energy over all degrees of freedom that enter quadratically in the Hamiltonian of the system. In particular for the average kinetic energy per degree of freedom, we have
\begin{equation}
  \langle \frac{1}{2} mv_{\alpha}^2 \rangle =   \frac{1}{2} k_BT
\end{equation}

In a simulation, we use this equation as an operational definition of the temperature. In practice, we would measure the total kinetic energy of the system and divide this by the number of degrees of freedom Nr (= 3N - 3 for a system of N particles with fixed total momentum1). As the total kinetic energy of a system fluctuates, so does the instantaneous temperature:

\begin{equation}
  T(t) = \sum_{i=1}^{N} \frac{m_iv_i^2(t)}{k_BN_f}
\end{equation}

Below is from the GROMACS manual:

While direct use of molecular dynamics gives rise to the NVE (constant number, constant volume, constant energy ensemble), most quantities that we wish to calculate are actually from a constant temperature (NVT) ensemble. 

To simulate a system under an ensemble which matches most experimental settings, where temperature and pressure is often held constant, we can modify the above introduced equations by adding additional terms which introduces a thermal heat bath (Nose-Hoover temperature coupling scheme) to which we couple our simulation to.
%\cite{[25] Nose ́, S. A molecular dynamics method for simulations in the canonical ensemble. Mol. Phys. 52:255–268, 1984.[26] Hoover, W. G. Canonical dynamics: equilibrium phase-space distributions. Phys. Rev. A 31:1695–1697, 1985.}

Using the same idea, we can couple the system to a pressure bath.  This method is called the Parrinello-Rahman pressure coupling, and the theory has been shown to correctly produce the NpT ensemble.
% \cite{Parrinello, M., Rahman, A. Polymorphic transitions in single crystals: A new molecular dynamics method. J. Appl. Phys. 52:7182–7190, 1981. [35] Nose ́, S., Klein, M. L. Constant pressure molecular dynamics for molecular systems. Mol. Phys. 50:1055–1076, 1983.}
% Should I include the equations? I haven't seen this in other thesis, perhaps its too much detail.

\subsection{Solvent Representation}
% Talk about how solvent is accounted for
The most realistic simulations treat the solvent environment by explicitly including the water molecules.  \textbf{How solvents are parameterized.  Can mention a few details about TIP3P water model. I actually don't know much about TIP3P} Solvent models are chosen to be consistent with the force field used for modelling the protein.  Although using explicit solvent leads to the most realistic simulations, these simulations are computationally expensive.  In some cases, instead of using explicit solvent molecules, implicit solvation can be used, that is the presence of solvent is approximated by using a potential. However, the implicit solvation suffers from several key weaknesses including a poor description of water-mediated interactions, directionality of water hydrogen bonds and only an approximate treatment of the nonpolar contribution to solvation.
Hence,  is much less accurate and can introduce force field inconsistencies leading to inaccurate results.\cite{Yang:2013ff}

\subsection{Periodic boundary conditions}
To treat a bulk system and avoid phase boundaries, periodic boundary conditions are used, where the the central simulation ``cell'' or ``box'' is replicated along each of it's X, Y and Z axis. Although this is unnatural, for sufficiently large systems, the errors induced by periodic boundaries are negligible. \textbf{Note I need a better explanation of PBC in my thesis as this was a topic of confusion in a couple of my committee meetings. Shouldn't be taken for granted. Also this could be a question during my defense as well.} REFS

\subsection{Ensembles}
Simulation systems are usually in the canonical ensemble (NVT) or  isothermo-isobaric (NPT) ensembles. These are statistical-mechanics ensembles used to relate the microscopic arrangements of the system to their thermodynamic properties such as the free energy. NVT is fixed number of particles and fixed volume and temperature.  NPT is fixed number of particles, constant pressure and temperature. This ensemble is the closest to experimental laboratory conditions, and hence simulations are often carried out in this fashion.

\section{Speeding up Molecular Dynamics Simulation}

\subsection{Better algorithms: Enhanced Sampling Techniques}

\subsubsection{Approaches for Determining Protein Conformational Ensembles}
It was mentioned in Chapter 1 that enhanced sampling may be used to speed up simulations.  Here we provide a high level overview of the approaches that were developed that have been useful in rational drug design.  These approaches were also widely used through out the thesis.

The use of temperature to overcome energetic barriers in the system. Here I can talk about what enhanced sampling is, and how they can be applied to understanding the properties of disordered peptides. Cite Sarah.  
I should read both Sarah's and Chris's thesis chapters relating to sampling.  Chris's is all about sampling. Since I refer extensively to enhanced sampling, I should make a mention of it.

\subsubsection{Approaches for Calculating Protein-Ligand Binding Affinities}
Note: This section could actually be put in the introduction.  As this may not require understanding the details of simulations.  Leave this hear for now and can move it as a section to the introduction ... maybe after Regis has seen both chapters together.

I can show the thermodynamic cycle for getting Delta G, like the one that I had in my transfer proposal.  Describe why this can be used to determine free energies.  And then mention the techniques that were invented in the past to do so. These approaches are not applicable to the problem of this thesis, but this background is provided to give the reader perspective. Point to another review article and say this is discussed in detail. I can discuss how to determine enthalpic and entropic contributions.


\subsection{Better Hardware}
Here can describe how better hardware - bigger and better machines have led to improvements and advancements in this field.

% Initial conditions -- initial positions of the atoms in the system

\section{Practical Aspects of Preparing a Biomolecular System for MD Simulation}

Coordinates for the initial input into a simulation software need some care to be prepared.  First the structure of protein or peptide to simulated need to be generated.  The starting structure of the target protein is usually X-ray crystal or NMR structure, which can be taken from the Protein Data Bank.  Then the simulation environment is prepared to closely mimic the experimental conditions (usually experimental crystallization conditions). For example, the pKa of a protein is typically predicted so that residues such as histidine are appropriately protonated.

For a simulation of a globular protein, after the protein structure has been prepared, cosolutes, salt and water molecules are added to create the final simulation box. However, for more complex simulation systems, such as those in the presence of a lipid bilayer, additional steps
% Also make sure that at this point everyone in my committee understands that there is no ionization in the simulations and no chemical reactions.  
 
At this stage,  an energy minimization algorithm is usually performed on the system to correct any bad steric contacts that might have been introduced when solvent was added. An example of a completed starting structure of a simulation is shown in Figure~\ref{fig:simulation_box}.

\begin{figure}
 \centering
 \includegraphics[width=5in]{figures/methods/simulation_box.pdf}
 \caption[An example of a MD simulation system]{An example of a MD simulation system.}
 \label{fig:simulation_box}
\end{figure}

Prior to performing production dynamics, which is what most of the results from a simulation will be based on, a short simulation to equilibrate the system is performed.  Equilibration is performed to allow the pressure and temperature of the system to stabilize. Here the stability of the protein is also validated to ensure that the protein structure has not grossly deviated (or unfolded) from the starting model.  

\begin{figure}
 \centering
 \includegraphics[width=5in]{figures/methods/equilibration.pdf}
 \caption[RMSD from the initial crystal structure calculated from an equilibriation]{Equilibration of the simulation system.}
 \label{fig:simulation_equilibration}
\end{figure}

\addcontentsline{toc}{section}{Bibliography}
\bibliographystyle{plain}
\bibliography{introduction}

%\section{Application of MD in structure-based drug discovery}
%
%% \1 (Why computational?) Can help us get protein dynamics is important for understanding protein function. We want to understand protein function because we want to be able to design drugs to cure diseases.
%
%% \1 A important application of MD simulation in biochemistry is the predicting of protein-ligand binding free energies.
%
%One application of MD simulations is in rational drug design. In recent years structure-based computer modeling of protein-ligand interactions have become a core component of modern drug discovery.  In early stage drug discovery, a target is identified along with putative binding sites.  Then, the structure of the target is determined using structural determination techniques such as NMR or X-ray crystallography.
%% [See Tom's thesis]
%Ligands which may act as potential drugs are expected to bind with a high affinity (low $K_d$) to the binding site. The goal is to discover,  high specificity inhibitors of a protein (usually an enzyme). In this process, the binding free energy of the ligand to its target is used to quantitatively evaluate how well a ligand binds. A crude estimate of the binding affinity can be obtained using computational docking methods, where the energetics of binding is typically estimated without accounting for either ligand or protein flexibility.  Although docking is fast, it is often inaccurate for identifying true drug candidates.
%
%With computer hardware becoming faster and cheaper, MD simulations can be used to rapidly prototype experimental ideas -- for example, one can perform computational alchemy, that is, ``mutate'' residues to test various hypotheses. Furthermore, simulations may be used to determine whether a chemical change will produce a more potent drug candidate. Currently, state of the art computational binding studies use MD simulations, where the protein and drug is allowed to relax and freely move about in the system. However, in the case of understanding a specific binding reaction often needed when developing an enzyme inhibitor, the ability to observe the relevant binding events is a low probability event on the timescale achievable by simulations. Therefore, a few enhanced techniques have been developed to accelerate this process.  They are briefly introduced below.
%
%\section{Free energy calculations}
%There are two advanced methods that have been developed for determining the absolute binding free energies in combination with MD simulations.
%
%% \3 Linear interaction energy -- Out of scope
%% \3 MM/PBSA - no explicit account for solvents -- Out of scope
%\subsection{Thermodynamic perturbation}
%This is the paper that I am taking most of the topics here\cite{Gilson:2007hz}
%
%\subsection{Thermodynamic integration}
%I actually have no idea what this is, and would be hard pressed to explain this properly.
%In this case, I don't think I should be covering these topics in my thesis.
%
%\subsection{Free energy perturbation}
%Alchemically change one molecule into another
%
%\section{Review of MD studies of amyloid inhibition by small molecules}
%% MD studies using brute-force sampling. Aid in medicinal chemistry by making suggestions for how to design new AD drugs.
%
%\begin{outline}
%	%  Excerpt from Transfer proposal
%	\1 In recent years, molecular dynamics simulations have been intensively used to investigate the molecular basis of the structure and stability of amyloid fibrils. 
%	
%	\1 MD simulations of Congo red binding have been done with the protofibril-like crystal structure composed of the segment GNNQQNY.\{Wu, 2007 \#621\}
%	
%	\1 A recent simulation study of an N-methylated peptide with A$\beta$16-22 models of amyloid aggregates has provided insight into the possible mechanism of action of peptide inhibitors of amyloid formation.\{Soto, 2007 \#597\} This peptide inhibitor was shown to preferentially bind monomers to form dimers, possibly acting to inhibit fibril formation by sequestering monomers. However, peptide-based inhibitors have poor pharmacological profiles as they are actively broken down by proteases in the stomach and are difficult to transport across the blood-brain barrier. In addition, these peptide inhibitors specifically target A$\beta$ and thus do not have the potential to treat multiple amyloid diseases.
%\end{outline}


% Scrap
% Motivate the use of MD simulations
% Describe the details of molecular dynamics simulations
% Review the basic derivations of MD simulation equations and why they work

% Molecular dynamics simulations are a useful tool to study the structure, dynamics, and interaction of biomolecules. 

% MD is a numerical algorithm which solves a system of Newtonian equations of motion, and provides as output the time-trajectory of atoms with femtosecond time resolution.

% To review before my defense

% Details of the mathematics (need to review the basic theory + Taylor series expansion) - get a book - tomorrow maybe?

% Why is MD correct? Describe the fundamental assumptions of MD. Here, I want to give the readers who aren't familiar with the methodological details of MD a sense of the rigorousness of MD.

% The assumption at a hand-wave level adapted from Tom's thesis

% - Relationship between force and energy 
% - Relationship between momentum and velocity 
% - Why numerical approach must be used (no analytical solution for N > 2)
% - How is the force field plugged into the general algorithm.






% \subsection{Setting up a MD simulation: practical aspects} - Here are some details to run a MD simulation of a biomolecular system.  Should I omit this from my introduction? This is not really essential in understanding the rest of my thesis, or is it? Perhaps this should go into an appendix instead -- this isn't interesting. The following steps are often used to setup and start a MD simulation system of a protein. First, a pre-determined structure, typically a coordinate structure from X-ray crystallography or NMR, or homology-modelling data. Then a force field and solvent is chosen.


%\section{Limitations of MD simulations}
% Rauscher:2010p5682,Rauscher:2009wr}  % Ref: DE Shaw and CN
% Schlick T (2010) Molecular modeling and simulation: an inter- disciplinary guide, interdisciplinary applied mathematics, vol 21, 2nd edn. Springer, New York


% Given the approximations of MD, MD is well-suited for probing the dynamics of complex systems on the timescale of picoseconds to seconds. ???

%Roughly speaking, one would like to run a simulation at least 10 times longer than the slowest important timescale in a system. Unfortunately, many biomolecular timescales exceed 1 ms, and in some cases by orders of magnitude (44).\cite{Zuckerman:2011dz} For molecular simulations to reliably predict, guide, and help explain experiment, these simulations require force fields of sufficient accuracy, adequate sampling of the relevant biomolecular motions (convergence) and a correct representation of the experimental conditions. Failures in any of these areas yield results which disagree with experiment.Until sampling is adequate, equilibrium properties computed from a simulation remain biased by the system�s starting state and no meaningful comparison with experiment is possible [6]. 

%% copied from Zukerman\cite{Zuckerman:2011dz}
%\textbf{Although routine explicit-solvent MD simula- tions are now four or five orders of magnitude longer (i.e., 100-103 ns currently), modern MD studies still appear to fall significantly short of what is needed for statistically valid equi- librium simulation (36, 38). Roughly speak- ing, one would like to run a simulation at least 10 times longer than the slowest impor- tant timescale in a system. Unfortunately, many biomolecular timescales exceed 1 ms, and in some cases by orders of magnitude (44).}
%% Other references that are relevant for sampling \cite{Grossfield:2009bn}
%
%% Copied from \cite{Mobley:2011ks}
%\textbf{For molecular simulations to reliably predict, guide, and help explain experiment, these simulations require force fields of sufficient accuracy, adequate sampling of the rel- evant biomolecular motions (convergence) and a correct representation of the experimental conditions. Failures in any of these areas yield results which disagree with experiment. 
%
%We may be tempted to blame disagreement with experiment on just one of these areas�force fields are perhaps the most common scapegoat, sometimes with good reason [1�5]�but any or all of the three may be a weak point. And, in some sense, adequate sampling is the weakest link. 
%
%Until sampling is adequate, equilibrium properties computed from a simulation remain biased by the system�s starting state and no meaningful comparison with experiment is possible [6]. 
%
%With an inadequate force field or a poor representation of the experimental conditions, results will disagree with experiment, but will be robust and improvement is relatively easy, but not so with inadequate sampling.}
%
%% copied from Sarah's thesis
%\textbf{Achieving complete (or even adequate) conformational sampling is one of the key challenges in biomolecular simulations.\cite{Gnanakaran:2003vh} The energy landscape of most biomolecules is �rugged� and the source of this ruggedness is two-fold. The energetic barriers separating accessible states are often larger than the available thermal energy, and there are typically a large number of states to be sampled. The timescales of many biomolecular processes, such as protein folding, are still far beyond the reach of our current computational capability, which is generally limited to the 10-8-10-7 s timescale for continuous simulations. For example, even the folding of small domains or secondary structure elements, such as ?-hairpins and mini-proteins, occur on the 1-10 ?s timescale.1 
%
%Consequently, conventional or �brute force� molecular dynamics (MD) alone is often insufficient to achieve complete Boltzmann sampling of the important states of many biologically relevant systems. For this reason, generalized-ensemble algorithms have become popular tools for conformational sampling.}

% Topics that I'm not going to discuss
%Evaluating the convergence of simulations is still a challenge - I don't want to  go into this as this is out of the scope of my thesis.
%Limitations in the accuracy of current force fields - Don't talk about this .. but this will probably come up during the defense regardless.